\section*{Overview}

Divide and conquer DP optimization is the first optimization I check when a transition looks like ``take the best split
point \(k\) for every position \(j\)'' and the straightforward implementation is one quadratic loop too slow.

\section*{When to Use It}

Use it when the recurrence has the shape
\[
dp[i][j] = \min_{k \le j} \bigl(dp[i-1][k-1] + C(k, j)\bigr),
\]
and the optimal split points satisfy:
\[
opt[i][j] \le opt[i][j+1].
\]

That monotonicity usually comes from a quadrangle inequality or a similar structure in the cost function.

\section*{Core Idea}

If the best split index for \(j\) never moves left when \(j\) increases, then after computing the optimum at the middle
position, the left half only needs to search smaller split indices and the right half only needs to search larger ones.

That is exactly the same pruning pattern as divide and conquer on a monotone answer boundary.

\section*{Key Insight}

The optimization does not come from algebraic magic. It comes from reducing the search interval for the best split at
each position. Once the opt indices are monotone, every recursive subproblem can ignore a large part of the \(k\)-range.

\section*{Operations / Main Technique}

\begin{itemize}
  \item verify or prove opt monotonicity,
  \item compute one DP row from the previous row,
  \item recurse on the left and right halves with narrowed split ranges.
\end{itemize}

\paragraph{Worked example.}
If the middle position \(j = 50\) is optimized by \(k = 18\), then every position left of \(50\) only needs to try
split points up to \(18\), while positions to the right only need to try split points from \(18\) onward.

\section*{Correctness Intuition}

The monotonicity of the optimal split points guarantees that recursion never removes the true optimum from a subproblem's
search interval. So every position still checks all of its valid candidates, just within a much smaller window.

\section*{Complexity Analysis}

For one row of length \(n\), the optimized computation is typically \(O(n \log n)\) if evaluating one candidate is
\(O(1)\). Over \(m\) rows, that becomes \(O(m n \log n)\).

\section*{Implementation}

The reference code is the reusable core: given a previous row, a current row, and a cost function, it fills one row via
the divide-and-conquer recursion.

\section*{Common Pitfalls}

\begin{itemize}
  \item Using the optimization without actually proving opt monotonicity.
  \item Off-by-one mistakes in the split range and in the meaning of \(dp[i-1][k-1]\).
  \item Forgetting that the cost function itself still needs to be \(O(1)\) or cheap enough.
\end{itemize}

\section*{Variants / Extensions}

\begin{itemize}
  \item Knuth optimization, which is even stronger but applies to a narrower family of recurrences.
  \item Convex hull trick when the transition can be rewritten as line queries.
  \item Aliens trick when the DP is wrapped inside a parametric search.
\end{itemize}

\section*{Practice Problems}

\begin{itemize}
  \item Partition DP with monotone optimal cut positions.
  \item Grouping or clustering DP where the segment cost is precomputable.
  \item Circular Barn and similar row-by-row partition problems.
\end{itemize}

\section*{References}

\begin{itemize}
  \item \href{https://usaco.guide/plat/DC-DP}{USACO Guide: Divide and Conquer DP}.
  \item \href{https://codeforces.com/catalog?locale=en}{Codeforces Catalog}.
  \item \href{https://cp-algorithms.com/}{cp-algorithms, via the references collected in USACO Guide and CP resources for this optimization family}.
\end{itemize}
